% !TeX program = lualatex
\documentclass[12pt,aspectratio=169]{beamer}
\usepackage[utf8]{inputenc}
\usepackage[T1]{fontenc}
\usepackage{amsmath,amsfonts,amssymb}
\usepackage{graphicx}
\usepackage{xcolor}
\usepackage{hyperref}
\usepackage{anyfontsize}
\usepackage{lmodern}
\usepackage{longtable}
\usepackage{array}
\usepackage{booktabs}

% ====== EXTREME FONT REDUCTION FOR MAX CONTENT ======
\renewcommand{\tiny}{\fontsize{3.5}{4.5}\selectfont}
\renewcommand{\scriptsize}{\fontsize{4.5}{5.5}\selectfont}
\renewcommand{\footnotesize}{\fontsize{5.5}{6.5}\selectfont}
\renewcommand{\small}{\fontsize{6.5}{7.5}\selectfont}
\renewcommand{\normalsize}{\fontsize{7.5}{8.5}\selectfont}
\renewcommand{\large}{\fontsize{8.5}{9.5}\selectfont}
\renewcommand{\Large}{\fontsize{9.5}{10.5}\selectfont}
\renewcommand{\LARGE}{\fontsize{10.5}{12}\selectfont}
\renewcommand{\huge}{\fontsize{11.5}{13.5}\selectfont}
\renewcommand{\Huge}{\fontsize{13}{16}\selectfont}

% ====== COLORS ======
\definecolor{redcorrect}{RGB}{200,30,30}
\definecolor{bluecorrect}{RGB}{30,80,200}
\definecolor{greencheck}{RGB}{0,150,50}
\definecolor{lightgray}{RGB}{245,245,245}
\definecolor{darkgray}{RGB}{40,40,40}
\definecolor{softyellow}{RGB}{255,248,200}
\definecolor{deepblue}{RGB}{20,60,150}
\definecolor{darkred}{RGB}{150,20,20}
\definecolor{orangewarning}{RGB}{255,140,0}
\definecolor{correctgreen}{RGB}{0,120,40}
\definecolor{trapcolor}{RGB}{180,80,0}
\definecolor{purpleconcept}{RGB}{120,50,180}
\definecolor{tealhard}{RGB}{0,120,120}

\setbeamercolor{title}{fg=white,bg=redcorrect}
\setbeamercolor{frametitle}{fg=white,bg=redcorrect}
\setbeamercolor{block title}{fg=white,bg=bluecorrect}
\setbeamercolor{block body}{fg=darkgray,bg=lightgray}
\setbeamercolor{normal text}{fg=darkgray}
\usecolortheme{default}

% ====== CUSTOM COMMANDS ======
\newcommand{\correct}[1]{\textcolor{greencheck}{\textbf{\checkmark}} #1}
\newcommand{\keypoint}[1]{\textcolor{deepblue}{\textbf{#1}}}
\newcommand{\hardconcept}[1]{\textcolor{darkred}{\textbf{#1}}}
\newcommand{\examtraps}[1]{\textcolor{trapcolor}{\textbf{⚠ TRAP:}} #1}
\newcommand{\recallbox}[1]{\fcolorbox{deepblue}{blue!5}{\parbox{0.88\textwidth}{\centering \textbf{REMEMBER:} #1}}}
\newcommand{\stepbox}[1]{%
	\begin{center}
		\fcolorbox{darkgray}{white}{%
			\parbox{0.92\textwidth}{\vspace{0.03cm} #1 \vspace{0.03cm}}%
		}
	\end{center}
}
\newcommand{\answerbox}[1]{\fcolorbox{correctgreen}{green!8}{\parbox{0.88\textwidth}{\centering \textcolor{correctgreen}{\textbf{\checkmark\ CORRECT:}} #1}}}
\newcommand{\wronganswer}[1]{\textcolor{redcorrect}{\textbf{✘}} #1}
\newcommand{\trapbox}[1]{\fcolorbox{trapcolor}{orange!8}{\parbox{0.88\textwidth}{\centering \textcolor{trapcolor}{\textbf{⚠ EXAM TRAP:}} #1}}}
\newcommand{\keyrule}[1]{\textcolor{tealhard}{\textbf{🔑}} #1}

\begin{document}
	
	% ================================================================
	% SLIDE 1: TITLE
	% ================================================================
	\begin{frame}
		\centering
		\vspace{0.5cm}
		{\fontsize{18}{22}\selectfont \textbf{CAMS Domain B}}\\[0.15cm]
		{\fontsize{14}{17}\selectfont \textbf{Global AFC Frameworks, Governance \& Regulations}}\\[0.2cm]
		{\fontsize{9}{11}\selectfont \textit{Advanced Explanations for Hard Topics \& Exam Traps}}\\[0.1cm]
		\rule{4cm}{0.5pt}\\[0.1cm]
		{\fontsize{7}{9}\selectfont FATF Mutual Evaluations · Immediate Outcomes · Wolfsberg · Basel Committee}\\[0.05cm]
		{\fontsize{7}{9}\selectfont EU Directives (3rd-6th) · AMLA · DORA · UK AML Framework}\\[0.05cm]
		{\fontsize{7}{9}\selectfont Three Lines of Defense · AFC Governance · Board Oversight · Risk Assessments}
		\vspace{0.3cm}
	\end{frame}
	
	% ================================================================
	% SLIDE 2: FATF - MUTUAL EVALUATIONS
	% ================================================================
	\begin{frame}
		\frametitle{\fontsize{11}{13}\selectfont \textbf{FATF Mutual Evaluations — Complete Breakdown}}
		\begin{block}{\fontsize{7}{9}\selectfont \textbf{What are Mutual Evaluations?}}
			\scriptsize
			\textbf{Peer reviews} of member countries' AML/CFT frameworks:
			\begin{itemize}
				\item Conducted by FATF and FATF-style regional bodies (FSRBs)
				\item Assess \textbf{both} technical compliance and effectiveness
				\item Lead to \textbf{public reports} identifying strengths and weaknesses
				\item Countries must \textbf{remediate} deficiencies
			\end{itemize}
		\end{block}
		\vspace{0.02cm}
		\begin{block}{\fontsize{7}{9}\selectfont \textbf{Technical Compliance vs Effectiveness — CRITICAL}}
			\scriptsize
			\begin{tabular}{|p{0.22\textwidth}|p{0.38\textwidth}|p{0.36\textwidth}|}
				\hline
				\keypoint{Factor} & \keypoint{Technical Compliance} & \keypoint{Effectiveness} \\
				\hline
				Focus & Does the law exist? & Is the law working? \\
				\hline
				Assessment & Rec 1-40 & 11 Immediate Outcomes \\
				\hline
				Rating & Compliant/LC/PC/NC & High/Substantial/Moderate/Low \\
				\hline
				Example & AML law passed & STRs used in convictions \\
				\hline
			\end{tabular}
		\end{block}
		\vspace{0.02cm}
		\begin{block}{\fontsize{7}{9}\selectfont \textbf{Effectiveness Assessment — What FATF Looks For}}
			\scriptsize
			\begin{itemize}
				\item \textbf{IO.3:} Supervisors use risk-based supervision and apply sanctions
				\item \textbf{IO.4:} FIs apply preventive measures consistent with ML/TF risks
				\item \textbf{IO.6:} STRs are used for ML/TF investigations
				\item \textbf{IO.7:} ML/TF investigations and prosecutions are effective
				\item \textbf{IO.8:} Confiscation of proceeds of crime is effective
			\end{itemize}
		\end{block}
		\vspace{0.02cm}
		\recallbox{\scriptsize \textbf{Key Rule:} Technical compliance ≠ Effectiveness. A country can have perfect laws but still be ineffective.}
	\end{frame}
	
	% ================================================================
	% SLIDE 3: FATF IMMEDIATE OUTCOMES (IOs) - DETAILED
	% ================================================================
	\begin{frame}
		\frametitle{\fontsize{11}{13}\selectfont \textbf{FATF 11 Immediate Outcomes — Complete Guide}}
		\begin{block}{\fontsize{7}{9}\selectfont \textbf{Immediate Outcomes — Effectiveness Assessment}}
			\scriptsize
			\begin{longtable}{|p{0.06\textwidth}|p{0.44\textwidth}|p{0.46\textwidth}|}
				\hline
				\keypoint{IO} & \keypoint{Title} & \keypoint{What It Assesses} \\
				\hline
				IO.1 & ML/TF Risks Understood & Country understands its ML/TF risks (NRA) \\
				\hline
				IO.2 & International Cooperation & Effective international info sharing \\
				\hline
				IO.3 & Supervision & Risk-based supervision; sanctions applied \\
				\hline
				IO.4 & Preventive Measures & FIs/DNFBPs apply AML/CFT measures \\
				\hline
				IO.5 & Legal Persons Not Misused & BO transparency; legal entities not used for ML \\
				\hline
				IO.6 & Financial Intelligence & STRs used for investigations \\
				\hline
				IO.7 & ML/TF Investigations & Effective investigations and prosecutions \\
				\hline
				IO.8 & Confiscation & Proceeds of crime are confiscated \\
				\hline
				IO.9 & TF Investigations & TF investigations and prosecutions \\
				\hline
				IO.10 & TF Preventive Measures & TF preventive measures are effective \\
				\hline
				IO.11 & PF Preventive Measures & Proliferation financing measures \\
				\hline
			\end{longtable}
		\end{block}
		\vspace{0.02cm}
		\begin{block}{\fontsize{7}{9}\selectfont \textbf{Supervisor Priorities — CAMS Exam Favorite}}
			\scriptsize
			Supervisors prioritize:
			\begin{itemize}
				\item \textbf{IO.3:} Financial supervision is risk-based and effective
				\item \textbf{IO.4:} FIs and DNFBPs apply preventive measures
			\end{itemize}
		\end{block}
		\vspace{0.02cm}
		\recallbox{\scriptsize \textbf{Key Rule:} IO.3 and IO.4 are \textbf{prioritized} by supervisors evaluating bank compliance.}
	\end{frame}
	
	% ================================================================
	% SLIDE 4: WOLFSBERG GROUP - PRIVATE SECTOR STANDARDS
	% ================================================================
	\begin{frame}
		\frametitle{\fontsize{11}{13}\selectfont \textbf{Wolfsberg Group — Private Sector Standards}}
		\begin{block}{\fontsize{7}{9}\selectfont \textbf{What is the Wolfsberg Group?}}
			\scriptsize
			\textbf{Association of 13 global banks} that develop standards for:
			\begin{itemize}
				\item AML/CFT
				\item KYC
				\item Correspondent banking
				\item Trade finance
				\item Private banking
				\item Risk assessment
			\end{itemize}
		\end{block}
		\vspace{0.02cm}
		\begin{block}{\fontsize{7}{9}\selectfont \textbf{Key Wolfsberg Documents — CAMS Exam}}
			\scriptsize
			\begin{itemize}
				\item \textbf{Wolfsberg AML Principles:} General AML framework for banks
				\item \textbf{Wolfsberg Correspondent Banking Due Diligence:} EDD for correspondent banking
				\item \textbf{Wolfsberg Private Banking Principles:} EDD for private banking clients
				\item \textbf{Wolfsberg Trade Finance Principles:} TBML risks and controls
				\item \textbf{Wolfsberg FAQs on AML Risk Assessment:} Risk assessment guidance
			\end{itemize}
		\end{block}
		\vspace{0.02cm}
		\begin{block}{\fontsize{7}{9}\selectfont \textbf{Correspondent Banking Requirements}}
			\scriptsize
			\begin{itemize}
				\item \textbf{EDD:} Enhanced due diligence on correspondent banks
				\item \textbf{Senior Approval:} Senior management approval required
				\item \textbf{Ongoing Monitoring:} Continuous monitoring of relationships
				\item \textbf{Information Sharing:} Share information on AML controls
				\item \textbf{Shell Banks:} \textbf{Absolutely prohibited} — no shell bank relationships
			\end{itemize}
		\end{block}
		\vspace{0.02cm}
		\recallbox{\scriptsize \textbf{Key Rule:} Wolfsberg standards are \textbf{not law} but are \textbf{industry best practices} that regulators expect.}
	\end{frame}
	
	% ================================================================
	% SLIDE 5: BASEL COMMITTEE - AML/CFT GUIDANCE
	% ================================================================
	\begin{frame}
		\frametitle{\fontsize{11}{13}\selectfont \textbf{Basel Committee — AML/CFT Guidance}}
		\begin{block}{\fontsize{7}{9}\selectfont \textbf{What is the Basel Committee?}}
			\scriptsize
			\textbf{Basel Committee on Banking Supervision} — sets global banking standards:
			\begin{itemize}
				\item BCBS 239: Risk Data Aggregation
				\item AML/CFT guidance for banks
				\item Consolidated supervision framework
				\item Sound management of operational risk
			\end{itemize}
		\end{block}
		\vspace{0.02cm}
		\begin{block}{\fontsize{7}{9}\selectfont \textbf{Three Lines of Defense — CRITICAL MODEL}}
			\scriptsize
			\begin{tabular}{|p{0.12\textwidth}|p{0.38\textwidth}|p{0.46\textwidth}|}
				\hline
				\keypoint{Line} & \keypoint{Role} & \keypoint{Example} \\
				\hline
				\textbf{1st Line} & Operational management & Business units identify, assess, control risks \\
				\hline
				\textbf{2nd Line} & Oversight & AML compliance, risk management, controls \\
				\hline
				\textbf{3rd Line} & Independent assurance & Internal audit \\
				\hline
			\end{tabular}
		\end{block}
		\vspace{0.02cm}
		\begin{block}{\fontsize{7}{9}\selectfont \textbf{Key Principles}}
			\scriptsize
			\begin{itemize}
				\item \textbf{Accountability:} First line owns and manages risk
				\item \textbf{Independence:} Second line must be independent from business
				\item \textbf{Objectivity:} Third line must be independent and objective
				\item \textbf{Separation of Duties:} Key roles must be separated
				\item \textbf{Board Oversight:} Board has ultimate responsibility
			\end{itemize}
		\end{block}
		\vspace{0.02cm}
		\recallbox{\scriptsize \textbf{Key Rule:} Three Lines of Defense = \textbf{operational management} (Line 1), \textbf{compliance oversight} (Line 2), \textbf{internal audit} (Line 3).}
	\end{frame}
	
	% ================================================================
	% SLIDE 6: EU AML DIRECTIVES - 3RD TO 6TH
	% ================================================================
	\begin{frame}
		\frametitle{\fontsize{11}{13}\selectfont \textbf{EU AML Directives — 3rd to 6th (Complete)}}
		\begin{block}{\fontsize{7}{9}\selectfont \textbf{3rd AML Directive (2005)}}
			\scriptsize
			\begin{itemize}
				\item Incorporated FATF 40 Recommendations
				\item Introduced \textbf{beneficial ownership} requirements
			 Required \textbf{risk-based approach}
				\item Expanded predicate offenses
				 Required group-wide AML programs
			\end{itemize}
		\end{block}
		\vspace{0.02cm}
		\begin{block}{\fontsize{7}{9}\selectfont \textbf{4th AML Directive (2015)}}
			\scriptsize
			\begin{itemize}
				\item Required \textbf{central registers} of beneficial owners
			 Enhanced customer due diligence requirements
			 Expanded PEP definition (domestic PEPs included)
				\item Required risk assessments at national level
				 Strengthened transaction monitoring requirements
			\end{itemize}
		\end{block}
		\vspace{0.02cm}
		\begin{block}{\fontsize{7}{9}\selectfont \textbf{5th AML Directive (2018)}}
			\scriptsize
			\begin{itemize}
				\item Extended to \textbf{virtual currencies} and \textbf{prepaid cards}
			 Enhanced transparency of BO registers (public access)
				\item Expanded DNFBP coverage
			 Strengthened customer due diligence
			Improved cooperation between FIUs
			\end{itemize}
		\end{block}
		\vspace{0.02cm}
		\begin{block}{\fontsize{7}{9}\selectfont \textbf{6th AML Directive (2020)}}
			\scriptsize
			\begin{itemize}
				\item \textbf{Harmonized predicate offenses} across member states
				\item \textbf{Harmonized criminal penalties}
				\item Extended liability to \textbf{legal persons}
				\item Strengthened cross-border cooperation
				\item Enhanced whistleblower protection
			\end{itemize}
		\end{block}
		\vspace{0.02cm}
		\recallbox{\scriptsize \textbf{Key Rule:} 3rd = BO; 4th = Central registers; 5th = Crypto/prepaid; 6th = Harmonized offenses \& penalties.}
	\end{frame}
	
	% ================================================================
	% SLIDE 7: EARLY EU DIRECTIVES - SHORTCOMINGS
	% ================================================================
	\begin{frame}
		\frametitle{\fontsize{11}{13}\selectfont \textbf{Early EU Directives (1st \& 2nd) — Shortcomings}}
		\begin{block}{\fontsize{7}{9}\selectfont \textbf{1st AML Directive (1991)}}
			\scriptsize
			\begin{itemize}
				\item Focused \textbf{only on drug trafficking}
				\item Lacked uniformity in implementation
				\item No requirement for suspicious transaction reporting
				\item No beneficial ownership requirements
				\item Limited scope (only credit institutions)
			\end{itemize}
		\end{block}
		\vspace{0.02cm}
		\begin{block}{\fontsize{7}{9}\selectfont \textbf{2nd AML Directive (2001)}}
			\scriptsize
			\begin{itemize}
				\item Expanded predicate offenses (but still limited)
				\item Still suffered from \textbf{divergent national interpretations}
				\item No risk-based approach
				\item No group-wide AML programs
				\item Weak transposition by member states
			\end{itemize}
		\end{block}
		\vspace{0.02cm}
		\begin{block}{\fontsize{7}{9}\selectfont \textbf{Major Shortcoming — CAMS Exam Favorite}}
			\scriptsize
			\begin{center}
				\textbf{Failure to ensure timely and complete transposition} by all member states.
			\end{center}
			\vspace{0.02cm}
			\scriptsize
			\begin{itemize}
				\item \textbf{Consequence:} \textbf{Regulatory arbitrage} — criminals exploited weak AML frameworks in certain EU countries.
				\item \textbf{Lesson:} Directives must be \textbf{transposed} into national law — and some countries delayed or did it poorly.
			\end{itemize}
		\end{block}
		\vspace{0.02cm}
		\recallbox{\scriptsize \textbf{Key Rule:} Early directives failed because of \textbf{poor transposition} and \textbf{regulatory arbitrage}.}
	\end{frame}
	
	% ================================================================
	% SLIDE 8: EU AMLA - ANTI-MONEY LAUNDERING AUTHORITY
	% ================================================================
	\begin{frame}
		\frametitle{\fontsize{11}{13}\selectfont \textbf{EU AMLA — Anti-Money Laundering Authority}}
		\begin{block}{\fontsize{7}{9}\selectfont \textbf{What is AMLA?}}
			\scriptsize
			\textbf{New EU authority} (established 2024) — centralizes AML/CFT supervision:
			\begin{itemize}
				\item Direct supervision of \textbf{high-risk cross-border} financial institutions
				\item Coordination of national AML supervisors
				\item Oversight of \textbf{selected} financial institutions
				\item Support for FIUs and law enforcement
			\end{itemize}
		\end{block}
		\vspace{0.02cm}
		\begin{block}{\fontsize{7}{9}\selectfont \textbf{AMLA's Role — CAMS Exam Favorite}}
			\scriptsize
			\begin{itemize}
				\item \textbf{Direct Supervision:} AMLA directly supervises selected \textbf{cross-border financial institutions} operating in multiple member states
				\item \textbf{Coordination:} Coordinates national supervisors to ensure consistent supervision
				\item \textbf{High-Risk Focus:} Focuses on institutions with \textbf{high-risk profiles}
				\item \textbf{Non-Compliance:} Can impose sanctions for non-compliance
			\end{itemize}
		\end{block}
		\vspace{0.02cm}
		\begin{block}{\fontsize{7}{9}\selectfont \textbf{When AMLA Assumes Oversight — CAMS Exam}}
			\scriptsize
			\textbf{Scenario:} EU-based fintech with operations in 5 member states flagged for non-compliance by multiple regulators.
			\par \textbf{Correct:} \textbf{AMLA} assumes oversight due to cross-border nature and risk profile.
		\end{block}
		\vspace{0.02cm}
		\recallbox{\scriptsize \textbf{Key Rule:} AMLA = \textbf{centralized supervision} for cross-border, high-risk EU financial institutions.}
	\end{frame}
	
	% ================================================================
	% SLIDE 9: DORA - REGULATORY EXPECTATIONS
	% ================================================================
	\begin{frame}
		\frametitle{\fontsize{11}{13}\selectfont \textbf{DORA — Digital Operational Resilience Act (Detailed)}}
		\begin{block}{\fontsize{7}{9}\selectfont \textbf{Scope — Who Must Comply?}}
			\scriptsize
			\begin{itemize}
				\item Banks and credit institutions
				\item Investment firms
				\item Crypto-asset service providers
				\item Payment institutions and e-money institutions
				\item Insurance and reinsurance companies
				\item Third-party ICT providers (critical ones)
			\end{itemize}
		\end{block}
		\vspace{0.02cm}
		\begin{block}{\fontsize{7}{9}\selectfont \textbf{Key Regulatory Expectation — CAMS Exam Favorite}}
			\scriptsize
			\begin{center}
				\textbf{Centralizing oversight} of ICT risk management across all member-state branches.
			\end{center}
		\end{block}
		\vspace{0.02cm}
		\begin{block}{\fontsize{7}{9}\selectfont \textbf{DORA Requirements — Complete List}}
			\scriptsize
			\begin{enumerate}
				\item \textbf{ICT Risk Governance Framework:} Unified across all EU operations
				\item \textbf{ICT Risk Assessments:} Identify and assess ICT risks
				\item \textbf{Incident Reporting:} Report ICT incidents to regulators
				\item \textbf{Third-Party ICT Provider Management:} Monitor and manage vendor risks
				\item \textbf{Resilience Testing:} Regular testing (including threat-led penetration testing)
				\item \textbf{Information Sharing:} Share threat intelligence with peers
			\end{enumerate}
		\end{block}
		\vspace{0.02cm}
		\trapbox{\scriptsize \textbf{Exam Trap:} \wronganswer{Relying solely on national regulators' cybersecurity requirements} is \textbf{insufficient}. DORA mandates centralized, pan-European oversight.}
	\end{frame}
	
	% ================================================================
	% SLIDE 10: UK AML FRAMEWORK - RISK-BASED APPROACH
	% ================================================================
	\begin{frame}
		\frametitle{\fontsize{11}{13}\selectfont \textbf{UK AML Framework — Complete Overview}}
		\begin{block}{\fontsize{7}{9}\selectfont \textbf{Key Legislation}}
			\scriptsize
			\begin{itemize}
				\item \textbf{Proceeds of Crime Act 2002 (POCA):} Criminalizes money laundering
				\item \textbf{Terrorism Act 2000:} Criminalizes terrorist financing
				\item \textbf{Money Laundering Regulations 2017:} AML/CFT obligations
				\item \textbf{Sanctions and Anti-Money Laundering Act 2018:} Sanctions framework
			\end{itemize}
		\end{block}
		\vspace{0.02cm}
		\begin{block}{\fontsize{7}{9}\selectfont \textbf{Risk-Based Approach — Primary Purpose — CAMS Exam}}
			\scriptsize
			\begin{center}
				\textbf{To allow firms to customize controls suited to their risk exposure.}
			\end{center}
		\end{block}
		\vspace{0.02cm}
		\begin{block}{\fontsize{7}{9}\selectfont \textbf{UK AML Supervisor — CAMS Exam Favorite}}
			\scriptsize
			\begin{itemize}
				\item \textbf{FCA (Financial Conduct Authority):} Main AML supervisor for financial services
				\item \textbf{HMT (HM Treasury):} Policy lead for AML/CFT
				\item \textbf{NCA (National Crime Agency):} UK FIU
				\item \textbf{JMLIT (Joint Money Laundering Intelligence Taskforce):} PPP between banks and law enforcement
			\end{itemize}
		\end{block}
		\vspace{0.02cm}
		\begin{block}{\fontsize{7}{9}\selectfont \textbf{UK-Specific Requirements}}
			\scriptsize
			\begin{itemize}
				\item \textbf{MLRO (Money Laundering Reporting Officer):} Required for all regulated firms
				\item \textbf{Consent Regime:} Must seek NCA consent before completing suspicious transactions
				\item \textbf{Beneficial Ownership Register:} Companies must file PSC (People with Significant Control) info
			\end{itemize}
		\end{block}
		\vspace{0.02cm}
		\recallbox{\scriptsize \textbf{Key Rule:} UK's risk-based approach = \textbf{customize controls to risk}, not eliminate controls.}
	\end{frame}
	
	% ================================================================
	% SLIDE 11: THREE LINES OF DEFENSE - DETAILED
	% ================================================================
	\begin{frame}
		\frametitle{\fontsize{11}{13}\selectfont \textbf{Three Lines of Defense — Complete Model}}
		\begin{block}{\fontsize{7}{9}\selectfont \textbf{First Line: Operational Management}}
			\scriptsize
			\begin{itemize}
				\item \textbf{Who:} Business units (e.g., retail banking, corporate banking, wealth management)
				\item \textbf{Responsibility:}
				\begin{itemize}
					\item Identify, assess, and control risks
				 Implement AML/CFT policies and procedures
					\item Train staff on AML/CFT requirements
				Escalate suspicious activity
				\end{itemize}
			\end{itemize}
		\end{block}
		\vspace{0.02cm}
		\begin{block}{\fontsize{7}{9}\selectfont \textbf{Second Line: Compliance Oversight}}
			\scriptsize
			\begin{itemize}
				\item \textbf{Who:} AML compliance, risk management, MLRO
			 \textbf{Responsibility:}
				\begin{itemize}
					\item Oversee and challenge first line
				Develop AML/CFT policies and procedures
					\item Conduct risk assessments
				 Monitor and report on compliance
					\item File SARs/STRs
				\end{itemize}
			\end{itemize}
		\end{block}
		\vspace{0.02cm}
		\begin{block}{\fontsize{7}{9}\selectfont \textbf{Third Line: Internal Audit}}
			\scriptsize
			\begin{itemize}
				\item \textbf{Who:} Internal audit function
			 \textbf{Responsibility:}
				\begin{itemize}
					\item Provide independent assurance
					\item Assess effectiveness of AML controls
					\item Report findings to board/audit committee
				 Follow up on remediation
				\end{itemize}
			\end{itemize}
		\end{block}
		\vspace{0.02cm}
		\recallbox{\scriptsize \textbf{Key Rule:} Lines must be \textbf{independent} but \textbf{collaborative}. First line owns risk; second line oversees; third line assures.}
	\end{frame}
	
	% ================================================================
	% SLIDE 12: AFC GOVERNANCE - BOARD OVERSIGHT
	% ================================================================
	\begin{frame}
		\frametitle{\fontsize{11}{13}\selectfont \textbf{AFC Governance — Board Oversight (Complete)}}
		\begin{block}{\fontsize{7}{9}\selectfont \textbf{Board Responsibilities for AFC}}
			\scriptsize
			\begin{itemize}
				\item \textbf{Ultimate Accountability:} Board is ultimately accountable for AML/CFT program
				\item \textbf{Approve Policies:} Approve AML/CFT policies and procedures
				\item \textbf{Approve MLRO:} Approve appointment of MLRO
				\item \textbf{Oversee Culture:} Ensure compliance culture is embedded
				\item \textbf{Challenge Management:} Challenge management on AFC performance
			\end{itemize}
		\end{block}
		\vspace{0.02cm}
		\begin{block}{\fontsize{7}{9}\selectfont \textbf{Board Oversight Mechanisms — CAMS Exam}}
			\scriptsize
			Which demonstrate effective board oversight? (Select Three.)
			\begin{enumerate}
				\item \textbf{Receiving regular reports on financial crime risks and control effectiveness}
				\item \textbf{Reviewing whether senior management is addressing audit findings}
				\item \textbf{Challenging management when AFC performance metrics show weakness}
			\end{enumerate}
		\end{block}
		\vspace{0.02cm}
		\begin{block}{\fontsize{7}{9}\selectfont \textbf{AFC Reports to Governance — CAMS Exam Favorite}}
			\scriptsize
			Appropriate contents for AFC reports to governance:
			\begin{enumerate}
				\item \textbf{Updates on regulatory interactions and supervisory reviews}
				\item \textbf{Assessment outcomes of control effectiveness on monitoring tools}
				\item \textbf{Investigation statuses of escalated SAR reviews}
			\end{enumerate}
		\end{block}
		\vspace{0.02cm}
		\trapbox{\scriptsize \textbf{Exam Trap:} \wronganswer{Participating in SAR investigations} — boards provide \textbf{oversight}, not operational involvement.}
	\end{frame}
	
	% ================================================================
	% SLIDE 13: ENTERPRISE-WIDE RISK ASSESSMENT (EWRA)
	% ================================================================
	\begin{frame}
		\frametitle{\fontsize{11}{13}\selectfont \textbf{EWRA — Enterprise-Wide Risk Assessment}}
		\begin{block}{\fontsize{7}{9}\selectfont \textbf{What is EWRA?}}
			\scriptsize
			A comprehensive, documented risk assessment that identifies, assesses, and mitigates ML/TF risks across the \textbf{entire organization}:
			\begin{itemize}
				\item All business units
				\item All product lines
				\item All customer segments
				\item All jurisdictions
				\item All distribution channels
			\end{itemize}
		\end{block}
		\vspace{0.02cm}
		\begin{block}{\fontsize{7}{9}\selectfont \textbf{EWRA Components — CAMS Exam Favorite}}
			\scriptsize
			\begin{itemize}
				\item \textbf{Inherent Risk Assessment:} Identify risks without controls
				\item \textbf{Control Effectiveness Assessment:} Assess how well controls mitigate risks
				\item \textbf{Residual Risk Assessment:} Remaining risk after controls
				\item \textbf{Risk Appetite Statement:} How much risk is acceptable?
				\item \textbf{Action Plan:} How to address residual risk?
			\end{itemize}
		\end{block}
		\vspace{0.02cm}
		\begin{block}{\fontsize{7}{9}\selectfont \textbf{When to Update EWRA — CAMS Exam}}
			\scriptsize
			EWRA must be updated when:
			\begin{itemize}
				\item \textbf{New products/services} are launched
				\item \textbf{New customer segments} are targeted
				\item \textbf{New jurisdictions} are entered
				\item \textbf{Significant business changes} occur
				\item \textbf{Regulatory changes} occur
				\item \textbf{New risk information} becomes available
			\end{itemize}
		\end{block}
		\vspace{0.02cm}
		\recallbox{\scriptsize \textbf{Key Rule:} EWRA must be \textbf{dynamic} — updated when new channels, products, or customers emerge.}
	\end{frame}
	
	% ================================================================
	% SLIDE 14: RISK-BASED APPROACH - FATF REC 1
	% ================================================================
	\begin{frame}
		\frametitle{\fontsize{11}{13}\selectfont \textbf{Risk-Based Approach — FATF Recommendation 1}}
		\begin{block}{\fontsize{7}{9}\selectfont \textbf{What is the Risk-Based Approach?}}
			\scriptsize
			\textbf{Foundational principle} of FATF framework:
			\begin{itemize}
				\item Countries and FIs must \textbf{identify, assess, and understand} ML/TF risks
				\item Apply \textbf{proportionate} measures to mitigate risks
				\item Allocate resources to \textbf{highest risk areas}
				\item \textbf{Not prescriptive} — flexible to specific risk profiles
			\end{itemize}
		\end{block}
		\vspace{0.02cm}
		\begin{block}{\fontsize{7}{9}\selectfont \textbf{Risk-Based Approach in Action — CAMS Exam}}
			\scriptsize
			\begin{itemize}
				\item \textbf{High Risk:} Enhanced due diligence (EDD), transaction monitoring, senior approval
				\item \textbf{Medium Risk:} Standard due diligence (CDD), regular monitoring
				\item \textbf{Low Risk:} Simplified due diligence (SDD), less frequent monitoring
			\end{itemize}
		\end{block}
		\vspace{0.02cm}
		\begin{block}{\fontsize{7}{9}\selectfont \textbf{What the Risk-Based Approach is NOT}}
			\scriptsize
			\begin{itemize}
				\item \textbf{Not} zero due diligence — some CDD always required
				\item \textbf{Not} de-risking — cannot simply decline entire categories
				\item \textbf{Not} one-size-fits-all — must be tailored to risk profile
				\item \textbf{Not} static — must be updated as risks change
			\end{itemize}
		\end{block}
		\vspace{0.02cm}
		\recallbox{\scriptsize \textbf{Key Rule:} The risk-based approach is about \textbf{proportionate controls}, not \textbf{no controls}.}
	\end{frame}
	
	% ================================================================
	% SLIDE 15: DNFBPS - FATF REC 22
	% ================================================================
	\begin{frame}
		\frametitle{\fontsize{11}{13}\selectfont \textbf{DNFBPs — FATF Recommendation 22 (Complete)}}
		\begin{block}{\fontsize{7}{9}\selectfont \textbf{What are DNFBPs?}}
			\scriptsize
			\textbf{Designated Non-Financial Businesses and Professions} — subject to AML/CFT obligations:
			\begin{itemize}
				\item Casinos (and online gambling)
				\item Real estate agents
				\item Lawyers, notaries, and other independent legal professionals
				\item Accountants and auditors
				\item Trust and company service providers (TCSPs)
				\item Dealers in precious metals and stones
			\end{itemize}
		\end{block}
		\vspace{0.02cm}
		\begin{block}{\fontsize{7}{9}\selectfont \textbf{DNFBP AML Obligations}}
			\scriptsize
			\begin{itemize}
				\item \textbf{CDD:} Identify and verify customers
				\item \textbf{BO Identification:} Identify beneficial owners
				\item \textbf{Record-Keeping:} Maintain records for 5+ years
				\item \textbf{STR Filing:} Report suspicious transactions
				\item \textbf{Internal Controls:} Implement AML programs
				\item \textbf{Training:} Train staff on AML/CFT
				\item \textbf{Supervision:} Subject to oversight by competent authority
			\end{itemize}
		\end{block}
		\vspace{0.02cm}
		\begin{block}{\fontsize{7}{9}\selectfont \textbf{DNFBP Risk Management — CAMS Exam Favorite}}
			\scriptsize
			Which strategies help manage DNFBP risks? (Select Two.)
			\begin{enumerate}
				\item \textbf{Creating a reputational risk or client selection committee}
				\item \textbf{Using standardized sector-specific questionnaires}
			\end{enumerate}
		\end{block}
		\vspace{0.02cm}
		\trapbox{\scriptsize \textbf{Exam Trap:} \wronganswer{Defaulting to EDD for all DNFBPs} is \textbf{not risk-based}. Apply proportionate measures.}
	\end{frame}
	
	% ================================================================
	% SLIDE 16: CORRESPONDENT BANKING - AML REQUIREMENTS
	% ================================================================
	\begin{frame}
		\frametitle{\fontsize{11}{13}\selectfont \textbf{Correspondent Banking — AML Requirements (Complete)}}
		\begin{block}{\fontsize{7}{9}\selectfont \textbf{What is Correspondent Banking?}}
			\scriptsize
			\textbf{Banking relationship} where one bank provides services to another:
			\begin{itemize}
				\item Payment clearing and settlement
				\item Cash management
				\item Trade finance
				\item Foreign exchange
			\end{itemize}
		\end{block}
		\vspace{0.02cm}
		\begin{block}{\fontsize{7}{9}\selectfont \textbf{Correspondent Banking AML Requirements — FATF Rec 13}}
			\scriptsize
			\begin{enumerate}
				\item \textbf{Gather Information:} Understand nature of respondent's business
				\item \textbf{Assess Controls:} Evaluate respondent's AML/CFT controls
				\item \textbf{Senior Approval:} Senior management approval required
				\item \textbf{Documentation:} Document the relationship
				\item \textbf{Ongoing Monitoring:} Monitor for changes in risk
				\item \textbf{EDD:} Enhanced due diligence for high-risk respondents
			\end{enumerate}
		\end{block}
		\vspace{0.02cm}
		\begin{block}{\fontsize{7}{9}\selectfont \textbf{Shell Banks — ABSOLUTE PROHIBITION}}
			\scriptsize
			\begin{center}
				\textbf{Shell banks are absolutely prohibited} — cannot establish or maintain correspondent relationships with shell banks.
			\end{center}
			\vspace{0.02cm}
			\scriptsize
			\begin{itemize}
				\item \textbf{Shell Bank:} Bank with no physical presence in any jurisdiction
				\item \textbf{Consequence:} Must \textbf{terminate} relationship if discovered
				\item \textbf{Wolfsberg:} Strongly reinforces prohibition
			\end{itemize}
		\end{block}
		\vspace{0.02cm}
		\recallbox{\scriptsize \textbf{Key Rule:} Shell banks are \textbf{absolutely prohibited} — no exceptions.}
	\end{frame}
	
	% ================================================================
	% SLIDE 17: IASB - INTERNATIONAL STANDARDS
	% ================================================================
	\begin{frame}
		\frametitle{\fontsize{11}{13}\selectfont \textbf{IASB — International Accounting Standards Board}}
		\begin{block}{\fontsize{7}{9}\selectfont \textbf{What is IASB?}}
			\scriptsize
			\textbf{International Accounting Standards Board} — sets global accounting standards:
			\begin{itemize}
				\item IFRS (International Financial Reporting Standards)
				\item IAS (International Accounting Standards)
				\item Applicable in 140+ countries
			\end{itemize}
		\end{block}
		\vspace{0.02cm}
		\begin{block}{\fontsize{7}{9}\selectfont \textbf{Relevance to AML — CAMS Exam}}
			\scriptsize
			\begin{itemize}
				\item \textbf{Financial Reporting:} Accurate financial statements required
				\item \textbf{Asset Valuation:} Valuation of assets for AML/CFT purposes
				\item \textbf{Disclosure:} Disclosure of off-balance sheet structures
				\item \textbf{Consolidation:} Group-wide financial reporting
				\item \textbf{Audit:} External audit of financial statements
			\end{itemize}
		\end{block}
		\vspace{0.02cm}
		\begin{block}{\fontsize{7}{9}\selectfont \textbf{Key IFRS Standards for AFC}}
			\scriptsize
			\begin{itemize}
				\item \textbf{IFRS 10:} Consolidated financial statements
				\item \textbf{IFRS 7:} Financial instruments disclosures
				\item \textbf{IFRS 9:} Financial instruments (impairment for losses)
				\item \textbf{IAS 24:} Related party disclosures
			\end{itemize}
		\end{block}
		\vspace{0.02cm}
		\recallbox{\scriptsize \textbf{Key Rule:} IASB standards ensure \textbf{transparency in financial reporting}, which supports AML efforts.}
	\end{frame}
	
	% ================================================================
	% SLIDE 18: IOSCO - SECURITIES REGULATION
	% ================================================================
	\begin{frame}
		\frametitle{\fontsize{11}{13}\selectfont \textbf{IOSCO — International Organization of Securities Commissions}}
		\begin{block}{\fontsize{7}{9}\selectfont \textbf{What is IOSCO?}}
			\scriptsize
			\textbf{Global standard-setter} for securities markets:
			\begin{itemize}
				\item 130+ member countries
				\item Sets international securities standards
				\item Promotes investor protection
				\item Facilitates cross-border cooperation
			\end{itemize}
		\end{block}
		\vspace{0.02cm}
		\begin{block}{\fontsize{7}{9}\selectfont \textbf{IOSCO Core Objectives — CAMS Exam Favorite}}
			\scriptsize
			Which support anti-financial crime efforts? (Select Two.)
			\begin{enumerate}
				\item \textbf{Enhancing investor protection through cooperation and enforcement}
				\item \textbf{Promoting financial stability by managing systemic risks and facilitating international information exchange}
			\end{enumerate}
		\end{block}
		\vspace{0.02cm}
		\begin{block}{\fontsize{7}{9}\selectfont \textbf{Relevance to AML — Complete}}
			\scriptsize
			\begin{itemize}
				\item \textbf{Market Abuse:} Combating insider trading and market manipulation
				\item \textbf{Fraud:} Preventing securities fraud
				\item \textbf{Information Sharing:} Cross-border enforcement cooperation
				\item \textbf{Disclosure:} Ensuring accurate disclosure of information
				\item \textbf{Investor Protection:} Protecting investors from financial crime
			\end{itemize}
		\end{block}
		\vspace{0.02cm}
		\recallbox{\scriptsize \textbf{Key Rule:} IOSCO objectives = \textbf{investor protection} and \textbf{financial stability} — both support AFC.}
	\end{frame}
	
	% ================================================================
	% SLIDE 19: G-20 ANTI-CORRUPTION WORKING GROUP
	% ================================================================
	\begin{frame}
		\frametitle{\fontsize{11}{13}\selectfont \textbf{G-20 Anti-Corruption Working Group}}
		\begin{block}{\fontsize{7}{9}\selectfont \textbf{What is the G-20 ACWG?}}
			\scriptsize
			\textbf{Working Group} under G-20 that:
			\begin{itemize}
				\item Coordinates international anti-corruption efforts
				\item Develops anti-corruption action plans
				\item Monitors implementation of commitments
				\item Promotes international cooperation
			\end{itemize}
		\end{block}
		\vspace{0.02cm}
		\begin{block}{\fontsize{7}{9}\selectfont \textbf{Key Focus Areas — CAMS Exam}}
			\scriptsize
			\begin{itemize}
				\item \textbf{Beneficial Ownership Transparency:} Improving BO disclosure
				\item \textbf{Asset Recovery:} Tracing and recovering stolen assets
				\item \textbf{Public Procurement:} Reducing corruption in procurement
				\item \textbf{International Cooperation:} Strengthening MLA and extradition
				\item \textbf{Whistleblower Protection:} Protecting those who report corruption
			\end{itemize}
		\end{block}
		\vspace{0.02cm}
		\begin{block}{\fontsize{7}{9}\selectfont \textbf{ACWG Contribution — CAMS Exam Favorite}}
			\scriptsize
			\textbf{Correct Answer:} Coordinating with member states to establish and implement \textbf{multi-year anti-corruption action plans.}
		\end{block}
		\vspace{0.02cm}
		\recallbox{\scriptsize \textbf{Key Rule:} G-20 ACWG = \textbf{coordination} and \textbf{action plans} for anti-corruption.}
	\end{frame}
	
	% ================================================================
	% SLIDE 20: IMF - AML/CFT EXPECTATIONS
	% ================================================================
	\begin{frame}
		\frametitle{\fontsize{11}{13}\selectfont \textbf{IMF AML/CFT Expectations — Complete}}
		\begin{block}{\fontsize{7}{9}\selectfont \textbf{IMF Role in AML/CFT}}
			\scriptsize
			\textbf{International Monetary Fund}:
			\begin{itemize}
				\item Assesses AML/CFT frameworks in member countries
				\item Provides technical assistance and capacity building
				\item Incorporates AML/CFT into FSAP (Financial Sector Assessment Programs)
				\item Links AML/CFT to financial stability
			\end{itemize}
		\end{block}
		\vspace{0.02cm}
		\begin{block}{\fontsize{7}{9}\selectfont \textbf{IMF Expectations — CAMS Exam Favorite}}
			\scriptsize
			\textbf{Correct Answer:} Reform AML/CFT frameworks and align them with \textbf{international expectations}.
		\end{block}
		\vspace{0.02cm}
		\begin{block}{\fontsize{7}{9}\selectfont \textbf{Important Concept — CAMS Exam}}
			\scriptsize
			\begin{center}
				\textbf{Improvements in economic measures can mask financial crime vulnerabilities.}
			\end{center}
			\vspace{0.02cm}
			\scriptsize
			\begin{itemize}
				\item Strong GDP growth \textbf{does not} mean strong AML controls
				\item Countries can have \textbf{economic success} but \textbf{weak AML frameworks}
				\item IMF assessment looks \textbf{beyond} economic indicators
			\end{itemize}
		\end{block}
		\vspace{0.02cm}
		\recallbox{\scriptsize \textbf{Key Rule:} Economic success ≠ AML effectiveness. IMF evaluates \textbf{financial crime vulnerabilities} separately.}
	\end{frame}
	
	% ================================================================
	% SLIDE 21: BASEL AML INDEX - COUNTRY RISK
	% ================================================================
	\begin{frame}
		\frametitle{\fontsize{11}{13}\selectfont \textbf{Basel AML Index — Country Risk Assessment}}
		\begin{block}{\fontsize{7}{9}\selectfont \textbf{What is the Basel AML Index?}}
			\scriptsize
			\textbf{Composite index} ranking countries based on AML/CFT risk:
			\begin{itemize}
				\item Published by Basel Institute on Governance
				\item Compiles vulnerability indicators across \textbf{legal, political, and financial} dimensions
				\item Used by FIs for \textbf{country risk assessment}
				\item Updated annually
			\end{itemize}
		\end{block}
		\vspace{0.02cm}
		\begin{block}{\fontsize{7}{9}\selectfont \textbf{Index Components — CAMS Exam Favorite}}
			\scriptsize
			\begin{itemize}
				\item \textbf{Quality of AML/CFT Framework:} Laws, regulations, enforcement
				\item \textbf{Corruption:} Level of corruption in the country
				\item \textbf{Financial Transparency:} BO disclosure, financial secrecy
				\item \textbf{Regulatory Capacity:} Effectiveness of supervision
				\item \textbf{Political Risk:} Stability, rule of law
				\item \textbf{Financial Crime:} Prevalence of financial crime
			\end{itemize}
		\end{block}
		\vspace{0.02cm}
		\begin{block}{\fontsize{7}{9}\selectfont \textbf{Why Use the Basel AML Index? — CAMS Exam}}
			\scriptsize
			\textbf{Correct Answer:} It compiles vulnerability indicators across \textbf{legal, political, and financial dimensions}.
		\end{block}
		\vspace{0.02cm}
		\recallbox{\scriptsize \textbf{Key Rule:} Basel AML Index = \textbf{country risk assessment} for AML/CFT due diligence.}
	\end{frame}
	
	% ================================================================
	% SLIDE 22: FINANCIAL SECRECY INDEX - DETAILED
	% ================================================================
	\begin{frame}
		\frametitle{\fontsize{11}{13}\selectfont \textbf{Financial Secrecy Index — Tax Justice Network}}
		\begin{block}{\fontsize{7}{9}\selectfont \textbf{What is the Financial Secrecy Index?}}
			\scriptsize
			Published by \textbf{Tax Justice Network}:
			\begin{itemize}
				\item Measures \textbf{financial secrecy} in jurisdictions
				\item Ranks countries on secrecy and scale of financial activity
				\item Used for \textbf{country risk assessment}
				\item Identifies \textbf{secrecy havens}
			\end{itemize}
		\end{block}
		\vspace{0.02cm}
		\begin{block}{\fontsize{7}{9}\selectfont \textbf{Factors Increasing FSI Rating — CAMS Exam Favorite}}
			\scriptsize
			Which contribute to elevated FSI rating? (Select Three.)
			\begin{enumerate}
				\item \textbf{Minimal reporting of beneficial ownership}
				\item \textbf{Extensive cross-border services to non-residents}
				\item \textbf{High ratio of secrecy score to global scale weight}
			\end{enumerate}
		\end{block}
		\vspace{0.02cm}
		\begin{block}{\fontsize{7}{9}\selectfont \textbf{FSI Implications}}
			\scriptsize
			\begin{itemize}
				\item \textbf{High FSI = High ML Risk:} Jurisdictions with high secrecy scores are higher risk
				\item \textbf{EDD Required:} Enhanced due diligence for clients with connections to high-FSI jurisdictions
				\item \textbf{Wolfsberg:} Cited in Wolfsberg Correspondent Banking Principles
			\end{itemize}
		\end{block}
		\vspace{0.02cm}
		\recallbox{\scriptsize \textbf{Key Rule:} High Financial Secrecy Index = \textbf{high money laundering risk}.}
	\end{frame}
	
	% ================================================================
	% SLIDE 23: TRANSPARENCY INTERNATIONAL - CPI
	% ================================================================
	\begin{frame}
		\frametitle{\fontsize{11}{13}\selectfont \textbf{Transparency International — Corruption Perceptions Index}}
		\begin{block}{\fontsize{7}{9}\selectfont \textbf{What is Transparency International?}}
			\scriptsize
			\textbf{Global NGO} fighting corruption:
			\begin{itemize}
				\item Publishes Corruption Perceptions Index (CPI)
				\item Advocates for anti-corruption reforms
				\item Conducts research on corruption
				\item Works with governments and businesses
			\end{itemize}
		\end{block}
		\vspace{0.02cm}
		\begin{block}{\fontsize{7}{9}\selectfont \textbf{Corruption Perceptions Index (CPI) — CAMS Exam}}
			\scriptsize
			\begin{itemize}
				\item \textbf{Ranks 180+ countries} by perceived corruption
				\item \textbf{Score:} 0 (highly corrupt) to 100 (very clean)
				\item \textbf{Source:} Expert surveys and business opinion
				\item \textbf{Use in AML:} Country risk assessment for AML/CFT due diligence
			\end{itemize}
		\end{block}
		\vspace{0.02cm}
		\begin{block}{\fontsize{7}{9}\selectfont \textbf{CPI and AML Risk}}
			\scriptsize
			\begin{itemize}
				\item \textbf{Low CPI Score = High Corruption Risk:} Countries with low CPI scores are higher risk
				\item \textbf{Correlation with ML:} Corruption is a \textbf{predicate offense} for money laundering
				\item \textbf{EDD:} Enhanced due diligence for clients with connections to low-CPI countries
			\end{itemize}
		\end{block}
		\vspace{0.02cm}
		\recallbox{\scriptsize \textbf{Key Rule:} CPI = \textbf{corruption risk}. Low CPI = high ML risk.}
	\end{frame}
	
	% ================================================================
	% SLIDE 24: PUBLIC-PRIVATE PARTNERSHIPS (PPP) - AML
	% ================================================================
	\begin{frame}
		\frametitle{\fontsize{11}{13}\selectfont \textbf{Public-Private Partnerships (PPP) — Complete Guide}}
		\begin{block}{\fontsize{7}{9}\selectfont \textbf{What are PPPs in AML?}}
			\scriptsize
			\textbf{Formal arrangements} between public authorities and private sector:
			\begin{itemize}
				\item FIUs and law enforcement
				\item Financial institutions (banks, fintechs)
				\item Information sharing on financial crime threats
			\end{itemize}
		\end{block}
		\vspace{0.02cm}
		\begin{block}{\fontsize{7}{9}\selectfont \textbf{PPP Principles — CAMS Exam Favorite}}
			\scriptsize
			\textbf{Correct Answer:} Collaboration allows \textbf{timely information sharing} and \textbf{joint targeting} of high-risk actors using shared intelligence.
		\end{block}
		\vspace{0.02cm}
		\begin{block}{\fontsize{7}{9}\selectfont \textbf{PPP Examples}}
			\scriptsize
			\begin{itemize}
				\item \textbf{UK JMLIT:} Joint Money Laundering Intelligence Taskforce
				\item \textbf{Australia Fintel:} Financial Intelligence Taskforce
				\item \textbf{Singapore ACIP:} Anti-Money Laundering Information Platform
				\item \textbf{US FinCEN Exchange:} Public-private information sharing
			\end{itemize}
		\end{block}
		\vspace{0.02cm}
		\begin{block}{\fontsize{7}{9}\selectfont \textbf{PPP Best Practices — CAMS Exam}}
			\scriptsize
			\textbf{Correct Approach:} Establishing a formal platform where regulators and institutions engage in \textbf{direct two-way communication} on specific threats.
		\end{block}
		\vspace{0.02cm}
		\recallbox{\scriptsize \textbf{Key Rule:} PPPs = \textbf{two-way information sharing} between public and private sectors.}
	\end{frame}
	
	% ================================================================
	% SLIDE 25: MUTUAL LEGAL ASSISTANCE TREATIES (MLATs)
	% ================================================================
	\begin{frame}
		\frametitle{\fontsize{11}{13}\selectfont \textbf{MLATs — Mutual Legal Assistance Treaties}}
		\begin{block}{\fontsize{7}{9}\selectfont \textbf{What are MLATs?}}
			\scriptsize
			\textbf{Bilateral/multilateral treaties} between countries:
			\begin{itemize}
				\item Facilitate cross-border legal cooperation
				\item Enforce criminal laws
				\item Share evidence and information
				\item Return fugitives and recovered assets
			\end{itemize}
		\end{block}
		\vspace{0.02cm}
		\begin{block}{\fontsize{7}{9}\selectfont \textbf{MLAT Objectives — CAMS Exam Favorite}}
			\scriptsize
			\textbf{Correct Answer:} Facilitating \textbf{structured cooperation} between foreign authorities for investigative assistance.
		\end{block}
		\vspace{0.02cm}
		\begin{block}{\fontsize{7}{9}\selectfont \textbf{MLAT Challenges — CAMS Exam Favorite}}
			\scriptsize
			\textbf{Scenario:} Country A has data protection rules preventing sharing of client identities. Country B bases investigations on BO disclosures.
			\par \textbf{Challenge:} Evidence exchanged under the MLAT will \textbf{not be admissible} in Country B's legal system.
		\end{block}
		\vspace{0.02cm}
		\begin{block}{\fontsize{7}{9}\selectfont \textbf{MLAT vs FATF Rec 36}}
			\scriptsize
			\begin{itemize}
				\item \textbf{FATF Rec 36:} Countries should provide the \textbf{widest possible range} of mutual legal assistance
				\item \textbf{Spontaneous Information Sharing:} Authorities may take \textbf{spontaneous measures} to share information
			\end{itemize}
		\end{block}
		\vspace{0.02cm}
		\recallbox{\scriptsize \textbf{Key Rule:} MLATs = \textbf{formal cooperation} for cross-border investigations.}
	\end{frame}
	
	% ================================================================
	% SLIDE 26: COOPERATION - FATF REC 40
	% ================================================================
	\begin{frame}
		\frametitle{\fontsize{11}{13}\selectfont \textbf{FATF Recommendation 40 — International Cooperation}}
		\begin{block}{\fontsize{7}{9}\selectfont \textbf{What is Rec 40?}}
			\scriptsize
			\textbf{International cooperation} between FIUs:
			\begin{itemize}
				\item FIU-to-FIU information sharing
				\item \textbf{Broadest possible range} of cooperation
				\item Spontaneous information sharing
				\item Exchange of financial intelligence
			\end{itemize}
		\end{block}
		\vspace{0.02cm}
		\begin{block}{\fontsize{7}{9}\selectfont \textbf{Rec 40 Requirements}}
			\scriptsize
			\begin{itemize}
				\item \textbf{Conditions:} Sharing must be done \textbf{securely} and \textbf{confidentially}
				\item \textbf{Spontaneous Sharing:} FIUs may share information \textbf{without request}
				\item \textbf{Reciprocity:} Should be possible on a \textbf{reciprocal} basis
				\item \textbf{Confidentiality:} Shared information must be kept \textbf{confidential}
				\item \textbf{Use:} Information can be used for \textbf{ML/TF investigations}
			\end{itemize}
		\end{block}
		\vspace{0.02cm}
		\begin{block}{\fontsize{7}{9}\selectfont \textbf{Rec 36 vs Rec 40 — CAMS Exam}}
			\scriptsize
			\begin{tabular}{|p{0.40\textwidth}|p{0.56\textwidth}|}
				\hline
				\keypoint{Rec 36} & \keypoint{Rec 40} \\
				\hline
				Mutual Legal Assistance (MLA) & FIU-to-FIU cooperation \\
				\hline
				Between countries (formal) & Between FIUs (operational) \\
				\hline
				Criminal investigations & Financial intelligence \\
				\hline
				Evidence sharing & Intelligence sharing \\
				\hline
			\end{tabular}
		\end{block}
		\vspace{0.02cm}
		\recallbox{\scriptsize \textbf{Key Rule:} Rec 36 = MLA (formal). Rec 40 = FIU-to-FIU (operational).}
	\end{frame}
	
	% ================================================================
	% SLIDE 27: FATF REC 21 - TIPPING OFF
	% ================================================================
	\begin{frame}
		\frametitle{\fontsize{11}{13}\selectfont \textbf{FATF Recommendation 21 — Tipping Off}}
		\begin{block}{\fontsize{7}{9}\selectfont \textbf{What is Tipping Off?}}
			\scriptsize
			\textbf{Disclosing to the subject} that:
			\begin{itemize}
				\item A suspicious transaction report (STR) has been filed
				\item They are under investigation
				\item Information has been shared with authorities
			\end{itemize}
		\end{block}
		\vspace{0.02cm}
		\begin{block}{\fontsize{7}{9}\selectfont \textbf{Why Prohibit Tipping Off?}}
			\scriptsize
			\begin{itemize}
				\item \textbf{Protects Investigations:} Preserves the integrity of the investigation
				\item \textbf{Prevents Evidence Destruction:} Subject may destroy evidence
				\item \textbf{Prevents Flight:} Subject may flee the jurisdiction
				\item \textbf{Protects Sources:} Protects the source of information
				\item \textbf{Prevents Witness Intimidation:} Subject may intimidate witnesses
			\end{itemize}
		\end{block}
		\vspace{0.02cm}
		\begin{block}{\fontsize{7}{9}\selectfont \textbf{Tipping Off Examples — CAMS Exam}}
			\scriptsize
			\begin{itemize}
				\item \textbf{DO NOT:} Tell the customer you filed an SAR
				\item \textbf{DO NOT:} Tell the customer they are under investigation
				\item \textbf{DO NOT:} Tell the customer information was shared with regulators
				\item \textbf{DO NOT:} Ask questions that would reveal the investigation
			\end{itemize}
		\end{block}
		\vspace{0.02cm}
		\begin{block}{\fontsize{7}{9}\selectfont \textbf{Exceptions}}
			\scriptsize
			\begin{itemize}
				\item \textbf{Legal Obligation:} When required by law (e.g., court order)
				\item \textbf{Law Enforcement:} When sharing with law enforcement
				\item \textbf{Regulators:} When sharing with supervisors
			\end{itemize}
		\end{block}
		\vspace{0.02cm}
		\recallbox{\scriptsize \textbf{Key Rule:} \textbf{NEVER} tip off the subject of an investigation.}
	\end{frame}
	
	% ================================================================
	% SLIDE 28: BENEFICIAL OWNERSHIP - FATF REC 24/25
	% ================================================================
	\begin{frame}
		\frametitle{\fontsize{11}{13}\selectfont \textbf{Beneficial Ownership — FATF Rec 24 \& 25}}
		\begin{block}{\fontsize{7}{9}\selectfont \textbf{What is Beneficial Ownership?}}
			\scriptsize
			\textbf{Natural person} who ultimately owns or controls:
			\begin{itemize}
				\item A legal person (company)
				\item A legal arrangement (trust)
			\end{itemize}
		\end{block}
		\vspace{0.02cm}
		\begin{block}{\fontsize{7}{9}\selectfont \textbf{Rec 24 — Legal Persons (Companies)}}
			\scriptsize
			\begin{itemize}
				\item Identify the \textbf{natural person} who ultimately \textbf{owns or controls} the company
				\item \textbf{Threshold:} Typically 25\% ownership (FATF)
				\item \textbf{Maintain BO Register:} Must maintain and update BO information
				\item \textbf{Verification:} Verify BO identity using reliable sources
			\end{itemize}
		\end{block}
		\vspace{0.02cm}
		\begin{block}{\fontsize{7}{9}\selectfont \textbf{Rec 25 — Legal Arrangements (Trusts)}}
			\scriptsize
			\begin{itemize}
				\item Identify: \textbf{Settlor}, \textbf{Trustee}, \textbf{Protector}, \textbf{Beneficiaries}
				\item Identify \textbf{natural persons} who ultimately \textbf{own or control}
				\item \textbf{Maintain BO Register:} Must maintain and update BO information
				\item \textbf{Verification:} Verify BO identity using reliable sources
			\end{itemize}
		\end{block}
		\vspace{0.02cm}
		\recallbox{\scriptsize \textbf{Key Rule:} BO = \textbf{natural person} who ultimately owns/controls. 25\% = typical threshold.}
	\end{frame}
	
	% ================================================================
	% SLIDE 29: RECORD-KEEPING - FATF REC 11
	% ================================================================
	\begin{frame}
		\frametitle{\fontsize{11}{13}\selectfont \textbf{Record-Keeping — FATF Recommendation 11}}
		\begin{block}{\fontsize{7}{9}\selectfont \textbf{Record-Keeping Requirements — Rec 11}}
			\scriptsize
			\begin{itemize}
				\item \textbf{Retention Period:} At least \textbf{5 years} after business relationship ends or transaction completed
				\item \textbf{Transaction Records:} Amount, currency, date, type of transaction, counterparty
				\item \textbf{Customer Records:} Copies of ID documents, account files, business correspondence
				\item \textbf{CDD Records:} Information obtained during CDD process
				\item \textbf{BO Records:} Beneficial ownership information
				\item \textbf{SAR/STR Records:} Suspicious transaction reports and supporting documentation
			\end{itemize}
		\end{block}
		\vspace{0.02cm}
		\begin{block}{\fontsize{7}{9}\selectfont \textbf{Why Record-Keeping Matters}}
			\scriptsize
			\begin{itemize}
				\item \textbf{Evidence:} Records are used as evidence in investigations
				\item \textbf{Audit Trail:} Provides audit trail for regulatory reviews
				\item \textbf{Reconstruction:} Allows reconstruction of transactions
				\item \textbf{Compliance:} Demonstrates compliance with AML/CFT obligations
				\item \textbf{Investigation:} Supports law enforcement investigations
			\end{itemize}
		\end{block}
		\vspace{0.02cm}
		\recallbox{\scriptsize \textbf{Key Rule:} Records must be kept for at least \textbf{5 years}.}
	\end{frame}
	
	% ================================================================
	% SLIDE 30: SAR/STR REPORTING - FATF REC 20
	% ================================================================
	\begin{frame}
		\frametitle{\fontsize{11}{13}\selectfont \textbf{SAR/STR Reporting — FATF Recommendation 20}}
		\begin{block}{\fontsize{7}{9}\selectfont \textbf{What is SAR/STR Reporting?}}
			\scriptsize
			\textbf{Suspicious Activity Report (US) / Suspicious Transaction Report (UK/EU)}:
			\begin{itemize}
				\item \textbf{SAR:} Suspicious Activity Report (US BSA)
				\item \textbf{STR:} Suspicious Transaction Report (UK/EU/Most countries)
				\item Filed with FIU or law enforcement
				\item Indicates \textbf{reasonable suspicion} of ML/TF
			\end{itemize}
		\end{block}
		\vspace{0.02cm}
		\begin{block}{\fontsize{7}{9}\selectfont \textbf{When to File — CAMS Exam Favorite}}
			\scriptsize
			\begin{itemize}
				\item \textbf{Reasonable Suspicion:} When there is reasonable suspicion of ML/TF
				\item \textbf{Attempted Transactions:} Even if the transaction was not completed
				\item \textbf{Structuring:} Attempts to avoid reporting thresholds
				\item \textbf{Inconsistent Information:} When customer information is inconsistent
				\item \textbf{No Required Minimum:} No minimum amount required (unlike CTRs)
			\end{itemize}
		\end{block}
		\vspace{0.02cm}
		\begin{block}{\fontsize{7}{9}\selectfont \textbf{US vs UK — CAMS Exam}}
			\scriptsize
			\begin{tabular}{|p{0.30\textwidth}|p{0.32\textwidth}|p{0.34\textwidth}|}
				\hline
				\keypoint{Factor} & \keypoint{US (SAR)} & \keypoint{UK (STR)} \\
				\hline
				Threshold & \$5,000+ (FIs) & No minimum \\
				\hline
				FIU & FinCEN & NCA \\
				\hline
				Consent Regime & Not required & Required for some \\
				\hline
				Protections & Safe harbor & Safe harbor \\
				\hline
			\end{tabular}
		\end{block}
		\vspace{0.02cm}
		\recallbox{\scriptsize \textbf{Key Rule:} SAR/STR = \textbf{reasonable suspicion} of ML/TF. No minimum amount required.}
	\end{frame}
	
	% ================================================================
	% SLIDE 31: TIPPING OFF - FATF REC 21 (DETAILED)
	% ================================================================
	\begin{frame}
		\frametitle{\fontsize{11}{13}\selectfont \textbf{Tipping Off — FATF Rec 21 (Detailed)}}
		\begin{block}{\fontsize{7}{9}\selectfont \textbf{Prohibition on Tipping Off}}
			\scriptsize
			\large
			\begin{center}
				\textbf{NEVER disclose SAR/STR investigation to the subject.}
			\end{center}
			\vspace{0.02cm}
			\scriptsize
			\begin{itemize}
				\item \textbf{Subject:} The person who is the subject of the SAR/STR
				\item \textbf{Third Parties:} Anyone who might inform the subject
				\item \textbf{Consequence:} Criminal penalties for tipping off
			\end{itemize}
		\end{block}
		\vspace{0.02cm}
		\begin{block}{\fontsize{7}{9}\selectfont \textbf{What Constitutes Tipping Off?}}
			\scriptsize
			\begin{itemize}
				\item \textbf{Direct Disclosure:} Telling the customer you filed an SAR
				\item \textbf{Indirect Disclosure:} Hinting that the customer is under investigation
				\item \textbf{Action:} Asking questions that would reveal the investigation
				\item \textbf{Behavior:} Changing treatment of the customer
				\item \textbf{Inaction:} Failing to follow normal procedures
			\end{itemize}
		\end{block}
		\vspace{0.02cm}
		\recallbox{\scriptsize \textbf{Key Rule:} Tipping off = \textbf{criminal offense} in most jurisdictions.}
	\end{frame}
	
	% ================================================================
	% SLIDE 32: FINAL EXAM TIPS - DOMAIN B
	% ================================================================
	\begin{frame}
		\frametitle{\fontsize{11}{13}\selectfont \textbf{Domain B — Critical Exam Tips}}
		\scriptsize
		\begin{block}{\fontsize{7}{9}\selectfont \textbf{Most Common Traps — Domain B}}
			\scriptsize
			\begin{itemize}
				\item \textbf{Technical Compliance vs Effectiveness:} Technical compliance = law exists. Effectiveness = law works. Both assessed.
				\item \textbf{IO.3 vs IO.4:} IO.3 = supervision. IO.4 = institutions applying measures.
				\item \textbf{Rec 36 vs Rec 40:} Rec 36 = MLA (formal). Rec 40 = FIU-to-FIU (operational).
				\item \textbf{3rd vs 4th vs 5th vs 6th:} 3rd = BO. 4th = registers. 5th = crypto/prepaid. 6th = harmonized.
				\item \textbf{AMLA vs National Regulators:} AMLA = cross-border. National = domestic.
				\item \textbf{DORA:} Centralized oversight of ICT risk — not just national cybersecurity.
				\item \textbf{Three Lines:} 1st = operations, 2nd = compliance, 3rd = audit.
				\item \textbf{UK vs US:} UK = STR, MLRO, Consent Regime. US = SAR, FinCEN.
				\item \textbf{Shell Banks:} Absolutely prohibited — no exceptions.
				\item \textbf{Tipping Off:} Never disclose SAR/STR investigation.
			\end{itemize}
		\end{block}
		\vspace{0.02cm}
		\begin{block}{\fontsize{7}{9}\selectfont \textbf{Exam Strategy — Domain B}}
			\scriptsize
			\begin{itemize}
				\item \textbf{Know the Frameworks:} FATF, Wolfsberg, Basel Committee, EU Directives
				\item \textbf{Know the Assessments:} Mutual evaluations assess technical compliance AND effectiveness
				\item \textbf{Know the Governance:} Three Lines of Defense, Board Oversight, EWRA
				\item \textbf{Know the Cooperation:} MLATs, PPPs, FIU-to-FIU, Rec 36/40
				\item \textbf{Think Governance:} When in doubt, it's about board oversight and risk-based supervision
			\end{itemize}
		\end{block}
		\vspace{0.02cm}
		\recallbox{\scriptsize \textbf{Final Rule:} Domain B = \textbf{global frameworks, governance, and cooperation}. Understand the structures.}
	\end{frame}
	
	% ================================================================
	% END SLIDE
	% ================================================================
	\begin{frame}
		\centering
		\vspace{1.5cm}
		{\fontsize{18}{22}\selectfont \textbf{Domain B — Complete}}\\[0.2cm]
		{\fontsize{11}{13}\selectfont Global AFC Frameworks, Governance \& Regulations}\\[0.1cm]
		\rule{4cm}{0.5pt}\\[0.1cm]
		{\fontsize{8}{10}\selectfont FATF Mutual Evaluations · 11 Immediate Outcomes · Wolfsberg}\\[0.05cm]
		{\fontsize{8}{10}\selectfont Basel Committee · EU Directives 3rd-6th · AMLA · DORA}\\[0.05cm]
		{\fontsize{8}{10}\selectfont UK AML Framework · Three Lines of Defense · Board Oversight}\\[0.05cm]
		{\fontsize{8}{10}\selectfont EWRA · Risk-Based Approach · DNFBPs · Correspondent Banking}\\[0.05cm]
		{\fontsize{8}{10}\selectfont MLATs · PPPs · FIU Cooperation · Rec 36/40 · Tipping Off}\\[0.05cm]
		{\fontsize{8}{10}\selectfont Beneficial Ownership · Record-Keeping · SAR/STR Reporting}\\[0.1cm]
		\rule{4cm}{0.5pt}\\[0.1cm]
		{\fontsize{8}{10}\selectfont \textit{All frameworks explained. All traps covered. Good luck on your CAMS exam!}}
		\vspace{0.5cm}
	\end{frame}
	
\end{document}